\documentclass{article}
\usepackage{amsmath,amssymb,amsthm}
\newtheorem{lemma}{Lemma}
\newtheorem{prop}{Proposition}
\begin{document}

\textbf{Subproblem~5 --- Obstruction from three edge directions (Solution~v11).}

\medskip
\noindent\textbf{Conventions and notation.}
We use the notation established in Subproblems~1--4.
\begin{itemize}
  \item $s_1,\dots,s_n$ are the vertices of $P=\operatorname{conv}(S)$ in cyclic
        (counterclockwise) order, indexed modulo~$n$.
  \item $a_i = s_{i+1}-s_i$, so $\sum_{i=1}^n a_i = 0$.
  \item For each $i$, there exists a direction $u_i$ in the interior of the normal cone
        $\mathcal{N}(s_i)$ of~$P$ at~$s_i$ such that $t_i\in T$ is the \emph{unique}
        maximiser of $\langle\cdot,u_i\rangle$ over~$T$.
        Since $u_i$ is in the interior of~$\mathcal{N}(s_i)$, the vertex~$s_i$ is
        also the \emph{unique} maximiser of $\langle\cdot,u_i\rangle$ over~$S$.
  \item $F_i = T\cap[t_i,t_{i+1}]=\{t_i+k\,a_i:0\le k\le m_i\}$ with
        $\sigma(t_i+k\,a_i)=(-1)^k\varepsilon_i$, $\varepsilon_i:=\sigma(t_i)$,
        and $t_{i+1}=t_i+m_i a_i$.
\end{itemize}
For $p\in T$ we have $\sigma_T(p)=\sigma(p)\in\{-1,+1\}$.
For $p\notin T$ we set $\sigma_T(p)=0$.

\bigskip
\noindent\textbf{Section~0: Structural consequences.}

\begin{lemma}[Lonely points]\label{lem:lonely}
$x_i:=s_i+t_i$ are pairwise distinct and
$\operatorname{supp}(\sigma)=\{x_1,\dots,x_n\}$.
\end{lemma}
\begin{proof}
Since $u_i$ uniquely maximises over $S\times T$, the pair $(s_i,t_i)$ is the unique maximiser,
so $\sigma(x_i)=\varepsilon_i\ne 0$ and the $x_i$ are distinct (if $x_i=x_j$ the pair
$(s_j,t_j)$ would give a second maximiser in direction $u_i$, contradicting $s_j\ne s_i$).
All $n$ points have $\sigma\ne 0$ and $|\operatorname{supp}(\sigma)|=n$, so they
exhaust the support.
\end{proof}

\begin{lemma}[Closure and parity]\label{lem:closure}
$\sum_{i}m_i a_i=0$ and $\sum_i m_i\equiv 0\pmod{2}$.
\end{lemma}
\begin{proof}
Telescoping: $\sum m_i a_i=t_{n+1}-t_1=0$.
The sign rule $\varepsilon_{i+1}=(-1)^{m_i}\varepsilon_i$ gives $(-1)^{\sum m_i}=1$.
\end{proof}

\begin{lemma}[Net sign]\label{lem:netsign}
$n\cdot\sum_{t\in T}\sigma(t)=\sum_i\varepsilon_i$.
\end{lemma}
\begin{proof}
$\sum_p\sigma(p)=n\sum_{t\in T}\sigma(t)$; also $\sum_p\sigma(p)=\sum_i\varepsilon_i$.
\end{proof}

\noindent\textbf{Non-lonely criterion.}
Since $\operatorname{supp}(\sigma)=\{x_1,\dots,x_n\}$:
\begin{equation}\label{eq:nonlonely}
  p\notin\{x_j\}
  \;\Longrightarrow\;
  \sigma(p)=0
  \;\text{i.e.,}\;
  \sum_{j=1}^n\sigma_T(p-s_j)=0.
\end{equation}
For $t\in T$ with $t\ne t_k$: $s_k+t\ne x_j$ for all $j$ (for $j=k$: $t\ne t_k$;
for $j\ne k$: $\langle s_k+t,u_j\rangle\le\langle s_k,u_j\rangle+\langle t_j,u_j\rangle
<\langle s_j,u_j\rangle+\langle t_j,u_j\rangle=\langle x_j,u_j\rangle$,
so $s_k+t\ne x_j$).
Therefore:
\begin{equation}\label{eq:sigma0}
  t\in T,\;t\ne t_k
  \;\Longrightarrow\;
  \sum_{j=1}^n\sigma_T(t+s_k-s_j)=0.
\end{equation}

\bigskip
\noindent\textbf{Section~1: The case $n=3$.}

Assume $n=3$; then $a_3=-a_1-a_2$ with $a_1,a_2$ linearly independent.
Define the linear functional $\beta:\mathbb{R}^2\to\mathbb{R}$ by $\beta(a_1)=0$, $\beta(a_2)=1$,
and set $h(p):=\beta(p-t_1)$ for points~$p$.
Then $h(t_1)=0$, $\beta(a_3)=-1$, $\beta(a_1)=0$, $\beta(a_2)=1$.

\begin{lemma}\label{lem:n3m}
$m_1=m_2=m_3=:m$, $m$ is even, and $m\ge 2$.
\end{lemma}
\begin{proof}
\textbf{Equality.}
From Lemma~\ref{lem:closure}: $(m_1-m_3)a_1+(m_2-m_3)a_2=0$.
Since $a_1,a_2$ are linearly independent: $m_1=m_2=m_3=:m$.
Parity: $3m\equiv 0\pmod 2$ forces $m$ even.

\textbf{$m\ge 2$.}
Suppose $m=0$.  Then $t_1=t_2=t_3=:t_0$.
Since $T$ has both signs, pick $t'\in T$ with $\sigma(t')\ne\sigma(t_0)$; in particular $t'\ne t_0$.

\textit{Case $h(t')=0$.}
Apply~\eqref{eq:sigma0} with $k=2$ at $t'$ (valid since $t'\ne t_2=t_0$):
\[
  \sigma_T(t'+a_1)+\sigma(t')+\sigma_T(t'-a_2)=0.
\]
Heights of the three terms: $0$, $0$, $h(t')-1=-1<0$.
So $\sigma_T(t'-a_2)=0$ (since $h(t'-a_2)=-1$; $h_{\min}\ge 0$ will be proved as Lemma~\ref{lem:betamin} for general $m\ge 0$, but for $m=0$ all face chains collapse to $t_0$ and any $T$-element at height $<0$ would descend indefinitely, as the same argument below shows).
Hence $\sigma_T(t'+a_1)=-\sigma(t')\ne 0$, so $t'+a_1\in T$ at height~$0$.
Applying~\eqref{eq:sigma0} with $k=1$ at $t'$:
$\sigma(t')+\sigma_T(t'-a_1)+\sigma_T(t'-a_1-a_2)=0$;
height of $t'-a_1-a_2$ is $-1<0$, so $\sigma_T(t'-a_1)=-\sigma(t')\ne 0$,
giving $t'-a_1\in T$ at height~$0$.
Iterating: $t'+ka_1\in T$ for all $k\in\mathbb{Z}$, a contradiction since $T$ is finite.

\textit{Case $h(t')>0$.}
Apply~\eqref{eq:sigma0} with $k=3$ at $t'$ (valid since $t'\ne t_3=t_0$):
\[
  \sigma_T(t'+a_1+a_2)+\sigma_T(t'+a_2)+\sigma(t')=0.
\]
So $\sigma_T(t'+a_1+a_2)+\sigma_T(t'+a_2)=-\sigma(t')\ne 0$.
At least one of $t'+a_2$, $t'+a_1+a_2$ is in $T$ at height $h(t')+1>h(t')$.
Repeating: $T$ contains elements at arbitrarily large heights, contradicting $T$ finite.

Therefore $m\ge 1$; since $m$ is even, $m\ge 2$.
\end{proof}

Set $\varepsilon:=\varepsilon_1$.  Since $m$ is even, $\varepsilon_2=\varepsilon_3=\varepsilon$.
Key heights: $h(t_1)=h(t_2)=0$, $h(t_3)=m$.

\begin{lemma}[$h_{\min}\ge 0$]\label{lem:betamin}
$h(t)\ge 0$ for every $t\in T$.
\end{lemma}
\begin{proof}
Suppose $h^*:=\min h(T)<0$, attained at $t^*$.
We verify $p:=s_3+(t^*-a_2)\notin\{x_j\}$ for $j=1,2,3$:
\begin{itemize}
\item $j=3$: $p=x_3=s_3+t_3$ iff $t^*-a_2=t_3$, i.e., $h(t^*)=h(t_3)+1=m+1>0>h^*$: not so.
\item $j=1$: $p=x_1=s_1+t_1$ iff $t^*-a_2=t_1+(s_1-s_3)=t_1-(a_1+a_2)$, i.e.,
  $t^*=t_1-a_1$, giving $h(t^*)=h(t_1)=0>h^*$: not so.
\item $j=2$: $p=x_2=s_2+t_2$ iff $t^*-a_2=t_2+(s_2-s_3)=t_2-a_2$, i.e., $t^*=t_2$,
  giving $h(t^*)=h(t_2)=0>h^*$: not so.
\end{itemize}
Hence $p\notin\{x_j\}$.  Apply~\eqref{eq:nonlonely} to $p$:
$\sigma_T(t^*+a_1)+\sigma_T(t^*)+\sigma_T(t^*-a_2)=0$.
(We used $s_3-s_1=a_1+a_2$, $s_3-s_2=a_2$, $s_3-s_3=0$.)
Since $h(t^*-a_2)=h^*-1<h^*=h_{\min}$: $\sigma_T(t^*-a_2)=0$.
Hence $\sigma_T(t^*+a_1)=-\sigma(t^*)\ne 0$; iterating (all iterates at height $h^*<0$ use the same equation) gives $t^*+ka_1\in T$ for all $k\ge 0$, contradiction.
\end{proof}

\begin{lemma}[$h_{\max}=m$, unique at $t_3$]\label{lem:betamax}
$h(t)\le m$ for all $t\in T$, and $t_3$ is the unique element with $h=m$.
\end{lemma}
\begin{proof}
\textbf{Upper bound.}
For $t\in T$ with $h(t)\ge m+1$ (so $t\ne t_3$): apply~\eqref{eq:sigma0} with $k=3$:
\[
  \sigma_T(t+a_1+a_2)+\sigma_T(t+a_2)+\sigma_T(t)=0.
\]
(Three terms, since $n=3$; $s_3-s_1=a_1+a_2$, $s_3-s_2=a_2$, $s_3-s_3=0$.)
Both $t+a_2$ and $t+a_1+a_2$ have height $\ge m+2$.
By induction, $T$ contains elements of arbitrarily large height: contradiction.

\textbf{Uniqueness at $h=m$.}
For $t\in T$ with $h(t)=m$, $t\ne t_3$: apply~\eqref{eq:sigma0} with $k=3$
(valid since $t\ne t_3$):
$\sigma_T(t+a_1+a_2)+\sigma_T(t+a_2)+\sigma(t)=0$.
Heights of $t+a_2,t+a_1+a_2$ are $m+1>m$, so both vanish.
Hence $\sigma(t)=0$, contradicting $t\in T$.
\end{proof}

\begin{lemma}[Staircase step]\label{lem:step}
For $t\in T$ with $0\le h(t)<m$ and $t\ne t_3$:
exactly one of $\{t+a_2,\,t+a_1+a_2\}$ lies in $T$ with sign $-\sigma(t)$;
the other has $\sigma_T=0$.
\end{lemma}
\begin{proof}
\textbf{Non-loneliness.}
$s_3+t=x_j$?
$j=3$: $t=t_3$, excluded.
$j=1$: $t=t_1+(s_1-s_3)=t_1-a_1-a_2$, so $h(t)=-1<0$: impossible by Lemma~\ref{lem:betamin}.
$j=2$: $t=t_2+(s_2-s_3)=t_2-a_2$, so $h(t)=h(t_2)-1=-1<0$: impossible.
Hence $s_3+t\notin\{x_j\}$.

\textbf{Equation.}
Apply~\eqref{eq:nonlonely} to $p=s_3+t$:
\[
  \sigma_T(t+a_1+a_2)+\sigma_T(t+a_2)+\sigma_T(t)=0.
\]
(Exactly three terms; $n=3$.)
Thus $\sigma_T(t+a_1+a_2)+\sigma_T(t+a_2)=-\sigma(t)\ne 0$.
Each term lies in $\{-1,0,+1\}$ and their sum has absolute value~$1$;
hence exactly one equals $-\sigma(t)$ and the other is~$0$.
\end{proof}

\begin{prop}[$n=3$ gives a contradiction]\label{prop:n3}
No finite signed set $T$ with both signs satisfies the hypotheses for $n=3$.
\end{prop}
\begin{proof}
By Lemma~\ref{lem:n3m}, $m\ge 2$ and $m$ is even.

\textbf{Starting element.}
$m$ is even and $m\ge 2$: $t_1+a_1\in F_1\subseteq T$ at height~$0$, sign~$-\varepsilon$.

\textbf{Staircase.}
Apply Lemma~\ref{lem:step} inductively starting from $t'_0:=t_1+a_1$ (sign~$-\varepsilon$):
for $k=0,1,\dots,m-1$,
starting from $t'_k$ at height~$k$ with sign $(-1)^k(-\varepsilon)$,
there is a unique $t'_{k+1}\in T$ at height $k+1$ with sign $(-1)^{k+1}(-\varepsilon)$.
(Uniqueness at $h=m$ from Lemma~\ref{lem:betamax} applies; all intermediate heights $\le m-1$.)

At $k=m-1$: $t'_m\in T$ with $h(t'_m)=m$, sign $(-1)^m(-\varepsilon)=-\varepsilon$ (since $m$ even).
By Lemma~\ref{lem:betamax}, $t'_m=t_3$.

\textbf{Sign contradiction.}
The actual sign of $t_3$ is $\varepsilon_3=(-1)^{m_1+m_2}\varepsilon=(-1)^{2m}\varepsilon=\varepsilon$ (since $m$ even).
The staircase assigns sign $-\varepsilon$ to $t_3$. Contradiction.
\end{proof}

\bigskip
\noindent\textbf{Section~2: The case $n=4$ with at least three edge-direction classes.}

From now until the end of Section~2, \textbf{$n=4$ and $P$ has $\ge 3$ edge-direction classes}.

\medskip
\noindent\textbf{Setup.}
Recall: direction classes are equivalence classes under positive scalar multiples.
A convex quadrilateral $P$ with $\ge 3$ direction classes is not a parallelogram
(a parallelogram has exactly 2 direction classes: $\{[a_1],[a_2]\}$ since $a_3\parallel a_1$ and $a_4\parallel a_2$).

Choose the labelling so that $a_1=s_2-s_1$ is the bottom edge and $\beta(s_j-s_1)>0$
for all $j\ge 3$ (possible since $P$ is convex; relabel cyclically if needed).
Set $a_2:=s_3-s_2$.
Define the linear functional $\beta:\mathbb{R}^2\to\mathbb{R}$ by $\beta(a_1)=0$, $\beta(a_2)=1$,
and $h(p):=\beta(p-t_1)$ for points~$p$.
Set $\mu:=\beta(a_3)$ and $c_j:=\beta(s_j-s_1)$ for $j=1,\dots,4$.
Then: $c_1=0$, $c_2=0$, $c_3=1$.

Since $n=4$: $a_4=-(a_1+a_2+a_3)$, so $\beta(a_4)=-(0+1+\mu)=-(1+\mu)$.
Key heights: $h(t_1)=h(t_2)=0$, $h(t_3)=m_2$, $h(t_4)=m_2+m_3\mu$.
The $\beta$-closure gives $m_2+m_3\mu+m_4\cdot(-(1+\mu))=0$, i.e.,
\begin{equation}\label{eq:closure4}
  m_2+m_3\mu = m_4(1+\mu).
\end{equation}
Bottom-edge property: $c_4=\beta(a_1+a_2+a_3)=1+\mu>0$, so $\mu>-1$.

The polygon has $\ge 3$ direction classes means $P$ is not a parallelogram.
The parallelogram condition ($[a_3]=[a_1]$ and $[a_4]=[a_2]$) requires
$a_3=\lambda a_1$ for some $\lambda>0$ (same direction class as $a_1$)
and $a_4=\nu a_2$ for some $\nu>0$ (same direction class as $a_2$).
If $a_3=\lambda a_1$: $\beta(a_3)=0$, i.e., $\mu=0$.
If additionally $a_4=\nu a_2$: $\beta(a_4)=-1$ gives $1+\mu=1$, consistent with $\mu=0$,
but then $a_4=-(a_1+a_2+a_3)=-(1+\lambda)a_1-a_2$ and requiring $a_4\parallel a_2$ forces
$1+\lambda=0$, i.e., $\lambda=-1$ (but direction class requires $\lambda>0$: contradiction).
Therefore, for $\ge 3$ direction classes, one of the following holds:
(A) $\mu\ne 0$ (i.e., $a_3$ is not a positive multiple of $a_1$); or
(B) $\mu=0$ but $a_3$ is not a positive multiple of $a_1$ (in fact $a_3=\lambda a_1$ with $\lambda<0$) and $a_4$ is not a positive multiple of $a_2$.

In both cases: the quadrilateral is not a parallelogram.

We split the proof into three sub-cases: $\mu\notin\mathbb{Z}$ (non-integer), $\mu=0$, and $\mu\in\mathbb{Z}_{>0}$.
Note $\mu\in\mathbb{Z}_{>0}$ with $\mu=0$ is separate from the truly non-integer case.

\begin{lemma}[$h_{\min}\ge 0$ for $n=4$]\label{lem:betamin4}
$h(t)\ge 0$ for every $t\in T$.
\end{lemma}
\begin{proof}
Suppose $h^*<0$, attained at $t'\in T$.  Since $h(t_2)=0>h^*$: $t'\ne t_2$.
Apply~\eqref{eq:nonlonely} to $p=s_2+t'$:
\[
  0=\sum_{j=1}^4\sigma_T(t'+(s_2-s_j))
  =\sigma_T(t'+a_1)+\sigma_T(t')+\sigma_T(t'-a_2)+\sigma_T(t'-a_2-a_3).
\]
Heights of terms: $h^*$, $h^*$, $h^*-1$, $h^*-1-\mu$.
Since $h^*-1<h^*$ and $h^*-1-\mu=h^*-(1+\mu)<h^*$ (using $1+\mu>0$),
both lower terms have height $<h^*=h_{\min}$, so they vanish.
Hence $\sigma_T(t'+a_1)=-\sigma(t')\ne 0$.
Applying~\eqref{eq:sigma0} with $k=1$ at $t'$ (heights $h^*$, $h^*$, $h^*-1$, $h^*-1-\mu$, lower terms vanish):
$\sigma(t')+\sigma_T(t'-a_1)=0\Rightarrow\sigma_T(t'-a_1)=-\sigma(t')\ne 0$.
Iterating: $t'+ka_1\in T$ for all $k\in\mathbb{Z}$, contradicting $T$ finite.
\end{proof}

\begin{lemma}[Height-zero structure]\label{lem:hzero}
$T\cap\{h=0\}=F_1=\{t_1+ka_1:0\le k\le m_1\}$.
\end{lemma}
\begin{proof}
Any $t\in T$ with $h(t)=0$ and $t\ne t_1,t_2$ satisfies: applying~\eqref{eq:sigma0}
with $k=2$ (valid since $t\ne t_2$):
$\sigma_T(t+a_1)+\sigma_T(t)+\sigma_T(t-a_2)+\sigma_T(t-a_2-a_3)=0$.
Heights: $0$,$0$,$-1$,$-1-\mu$.  Both $-1<0$ and $-1-\mu<0$ (since $\mu>-1$):
lower terms vanish.  Thus $\sigma_T(t+a_1)=-\sigma(t)\ne 0$.
Similarly from $k=1$: $\sigma_T(t-a_1)=-\sigma(t)\ne 0$.
Hence $\{t+ka_1:k\in\mathbb{Z}\}\subseteq T$, contradicting $T$ finite.

Therefore every $t\in T\cap\{h=0\}$ satisfies $t=t_1$ or $t=t_2$ (the excluded points)
or is one of $t_1+a_1,\dots,t_1+(m_1-1)a_1$ (interior of $F_1$, also excluded from $k=1$ or $k=2$ separately but handled similarly via the sign recurrence within $F_1$).
More precisely, the elements of $F_1$ themselves satisfy all equations consistently
(each satisfies $\sigma_T(t+a_1)=-\sigma(t)$ or $t+a_1=t_{i+1}$),
and any height-$0$ element outside $F_1$ would start a bi-infinite chain.
Thus $T\cap\{h=0\}=F_1$.
\end{proof}

\begin{lemma}[Height maximum via $s_4$-equation]\label{lem:betamax4mu}
For $\mu>0$: $h_{\max}=h(t_4)=m_2+m_3\mu$, achieved uniquely at $t_4$.
\end{lemma}
\begin{proof}
Let $t_*\in T$ with $h(t_*)=h_{\max}$ and $t_*\ne t_4$.
Apply~\eqref{eq:sigma0} with $k=4$ at $t_*$:
\[
  \sigma_T(t_*+a_1+a_2+a_3)+\sigma_T(t_*+a_2+a_3)+\sigma_T(t_*+a_3)+\sigma_T(t_*)=0.
\]
Heights of the first three terms: $h_{\max}+(1+\mu)$, $h_{\max}+(1+\mu)$, $h_{\max}+\mu$.
Since $\mu>0$: all three heights exceed $h_{\max}$, so all three $\sigma_T$ values are~$0$.
Hence $\sigma_T(t_*)=0$, contradicting $t_*\in T$.
Therefore $h_{\max}$ is achieved only at $t_4$, and $h_{\max}=m_2+m_3\mu$.
\end{proof}

\begin{lemma}[Height maximum for $\mu<0$]\label{lem:betamax4neg}
For $\mu<0$: $h_{\max}=m_2$, achieved uniquely at $t_3$.
\end{lemma}
\begin{proof}
Let $t_*\in T$ with $h(t_*)=h_{\max}$ and $t_*\ne t_3$.
Apply~\eqref{eq:sigma0} with $k=3$ at $t_*$:
\[
  \sigma_T(t_*+a_1+a_2)+\sigma_T(t_*+a_2)+\sigma_T(t_*)+\sigma_T(t_*-a_3)=0.
\]
Heights: $h_{\max}+1$, $h_{\max}+1$, $h_{\max}$, $h_{\max}-\mu=h_{\max}+|\mu|$.
Since $\mu<0$: $|\mu|>0$, so $h_{\max}+|\mu|>h_{\max}$, and $h_{\max}+1>h_{\max}$.
All three terms other than $\sigma_T(t_*)$ are at heights exceeding $h_{\max}$, so they vanish.
Hence $\sigma_T(t_*)=0$, contradicting $t_*\in T$.
Therefore $h_{\max}$ is achieved only at $t_3$, and $h_{\max}=m_2$.
\end{proof}

\begin{lemma}[Height maximum for $\mu=0$]\label{lem:betamax4zero}
For $\mu=0$: $h_{\max}=m_2$.
\end{lemma}
\begin{proof}
Suppose $t\in T$ with $h(t)>m_2$.  Then $t\ne t_3,t_4$ (since $h(t_3)=h(t_4)=m_2$).
Apply~\eqref{eq:sigma0} with $k=4$ at $t$ (valid since $t\ne t_4$):
\[
  \sigma_T(t+a_1+a_2+a_3)+\sigma_T(t+a_2+a_3)+\sigma_T(t+a_3)+\sigma_T(t)=0.
\]
Heights: $h(t)+1+\mu=h(t)+1$, $h(t)+1$, $h(t)$, $h(t)$.
Both terms at height $h(t)+1>h(t)$ vanish if $h_{\max}=h(t)$;
so $\sigma_T(t+a_3)+\sigma(t)=0$, giving $\sigma_T(t+a_3)=-\sigma(t)\ne 0$.
Since $\beta(a_3)=\mu=0$: $h(t+a_3)=h(t)$, so $t+a_3\in T$ at the same height.
If $t+a_3\ne t_4$: apply the same argument again; iterating gives $t+ka_3\in T$ for all $k\ge 0$,
all at height $h(t)>m_2$.  Since $T$ is finite, this is a contradiction.
(If $t+a_3=t_4$: then $t=t_4-a_3$, but $h(t_4-a_3)=h(t_4)-\mu=m_2$, contradicting $h(t)>m_2$.)
Hence no $T$-element has height $>m_2$, so $h_{\max}=m_2$.
\end{proof}

\begin{lemma}[No elements at non-integer heights]\label{lem:nogaps}
For $\mu\notin\mathbb{Z}$\textup{:}
$T$ has no element at any height $h\in\mathbb{R}\setminus\mathbb{Z}$ with $0<h<m_2$.
\end{lemma}
\begin{proof}
Since $T$ is finite, if any such element exists there is one at the minimum
such height.  Let $h_0=\min\{h(t):t\in T,\,h(t)\in(0,m_2)\setminus\mathbb{Z}\}$,
and let $t'\in T$ with $h(t')=h_0$.

Since $h_0>0$: $h(t')>0=h(t_1)=h(t_2)$, so $t'\ne t_1$ and $t'\ne t_2$.

Apply~\eqref{eq:sigma0} with $k=1$ at $t'$:
\begin{equation}\label{eq:k1t}
  \sigma_T(t')+\sigma_T(t'-a_1)+\sigma_T(t'-a_1-a_2)+\sigma_T(t'-a_1-a_2-a_3)=0,
\end{equation}
with term heights $h_0,\,h_0,\,h_0-1,\,h_0-1-\mu$.

Apply~\eqref{eq:sigma0} with $k=2$ at $t'$:
\begin{equation}\label{eq:k2t}
  \sigma_T(t'+a_1)+\sigma_T(t')+\sigma_T(t'-a_2)+\sigma_T(t'-a_2-a_3)=0,
\end{equation}
with term heights $h_0,\,h_0,\,h_0-1,\,h_0-1-\mu$.

\textbf{Lower-term analysis.}
The term at height $h_0-1$: since $h_0$ is non-integer and $1$ is an integer,
$h_0-1$ is non-integer.  If $h_0-1<0$: $\sigma_T=0$ by Lemma~\ref{lem:betamin4}.
If $0<h_0-1<m_2$: then $h_0-1<h_0$ is a smaller non-integer height in $(0,m_2)$,
so $T$ has no element there by minimality of $h_0$.  Thus $\sigma_T(t'-a_2)=\sigma_T(t'-a_1-a_2)=0$.

The term at height $h_0-1-\mu$: we have $h_0-1-\mu = h_0-(1+\mu)$ and
$1+\mu>0$ (since $\mu>-1$), so $h_0-1-\mu<h_0$.  Two sub-cases:

\textbf{Sub-case A:} $h_0-1-\mu\notin\mathbb{Z}$.  Then $h_0-1-\mu$ is a non-integer.
If $h_0-1-\mu\le 0$: $\sigma_T=0$ (Lemma~\ref{lem:betamin4} or $\le 0$ directly).
If $0<h_0-1-\mu<m_2$: $h_0-1-\mu<h_0$ is a smaller non-integer in $(0,m_2)$,
so $T$ has no element there by minimality.
If $h_0-1-\mu\ge m_2$: since $h_0<m_2$ and $1+\mu>0$, we have $h_0-1-\mu=h_0-(1+\mu)<m_2-(1+\mu)$.
For $\mu<0$: $h_0-1-\mu=h_0+|\mu|-1$, and $h_0<m_2$ with $|\mu|<1$ (since $\mu>-1$) gives
$h_0+|\mu|-1<m_2+|\mu|-1<m_2$ (since $|\mu|<1$): contradiction with $h_0-1-\mu\ge m_2$.
For $\mu>0$: $h_0-1-\mu<h_0-1<m_2$: the case $\ge m_2$ is impossible.
Hence $\sigma_T(t'-a_2-a_3)=\sigma_T(t'-a_1-a_2-a_3)=0$ in Sub-case~A.

Subcase A conclusion: From~\eqref{eq:k1t}: $\sigma(t')+\sigma_T(t'-a_1)=0$.
From~\eqref{eq:k2t}: $\sigma_T(t'+a_1)+\sigma(t')=0$.
So $\sigma_T(t'-a_1)=\sigma_T(t'+a_1)=-\sigma(t')\ne 0$: a bi-infinite chain $\{t'+ka_1:k\in\mathbb{Z}\}\subseteq T$, contradicting $T$ finite.

\textbf{Sub-case B:} $h_0-1-\mu\in\mathbb{Z}$.  Write $j:=h_0-1-\mu\in\mathbb{Z}$.
Since $\mu\notin\mathbb{Z}$ and $h_0\notin\mathbb{Z}$: $j=h_0-1-\mu$ is well-defined with $j\ge 0$
(since $h_0-1-\mu=h_0-(1+\mu)$ and $h_0>0$ while $1+\mu>0$; also $h_0\ge 1+\mu$ for $j\ge 0$
which holds since $h_0\in(0,m_2)$ and $1+\mu>0$, so $j$ can be $0$ if $h_0=1+\mu$).

Let $\gamma:=\sigma_T(t'-a_2-a_3)\in\{-1,0,+1\}$ (the height-$j$ term).
Let $\gamma':=\sigma_T(t'-a_1-a_2-a_3)\in\{-1,0,+1\}$ (also at height $j$).
From~\eqref{eq:k2t}: $\sigma_T(t'+a_1)+\sigma(t')=-\gamma$.
From~\eqref{eq:k1t}: $\sigma(t')+\sigma_T(t'-a_1)=-\gamma'$.

Both $\sigma_T(t'+a_1)$ and $\sigma_T(t'-a_1)$ must lie in $\{-1,0,+1\}$.
$\sigma(t')=-\sigma(t')\cdot(-1)\in\{-1,+1\}$.

If $\gamma=\sigma(t')$ (same sign): $\sigma_T(t'+a_1)=-\sigma(t')-\gamma=-2\sigma(t')\notin\{-1,0,+1\}$: \textbf{impossible}, contradiction.

If $\gamma=-\sigma(t')$: $\sigma_T(t'+a_1)=0$.
If additionally $\gamma'=\sigma(t')$: $\sigma_T(t'-a_1)=-\sigma(t')-\gamma'=-2\sigma(t')$: impossible, contradiction.
If $\gamma'=-\sigma(t')$: $\sigma_T(t'-a_1)=0$.

Now $j=h_0-1-\mu\ge 0$.  If $j=0$: height-$0$ terms $\gamma,\gamma'$ come from $T\cap\{h=0\}=F_1$ (Lemma~\ref{lem:hzero}):
$\gamma=\sigma_T(t'-a_2-a_3)$, element at height $0$.  If this element is in $F_1$: $\gamma=(-1)^k\varepsilon_1\ne 0$ for some $k$.
Then $\gamma$ and $\sigma(t')$ either agree (contradiction as above) or disagree.
If they disagree: $\sigma_T(t'+a_1)=0$.  Consider $\gamma'=\sigma_T(t'-a_1-a_2-a_3)$ at height $0$.
If $\gamma'=(-1)^{k'}\varepsilon_1$ (element of $F_1$): $\gamma'$ and $\sigma(t')$ either agree ($\sigma_T(t'-a_1)$ magnitude 2, contradiction) or disagree ($\sigma_T(t'-a_1)=0$).
In the surviving sub-case $\sigma_T(t'+a_1)=\sigma_T(t'-a_1)=0$:
apply~\eqref{eq:sigma0} with $k=1$ at $t'-a_1$ (height $h_0$, $\ne t_1$):
$\sigma_T(t'-a_1)+\sigma_T(t'-2a_1)+\sigma_T(t'-2a_1-a_2)+\sigma_T(t'-2a_1-a_2-a_3)=0$.
Heights $h_0$,$h_0$,$h_0-1$,$j=0$.
Lower term $\sigma_T(t'-2a_1-a_2-a_3)$ at height $0$: element of $F_1$ or $0$.
Since $\sigma_T(t'-a_1)=0$ (not in $T$, so can't apply $k=1$ at $t'-a_1$; instead note: since
$\sigma_T(t'+a_1)=\sigma_T(t'-a_1)=0$ and $t'$ is isolated at height $h_0$, the only way
$\sigma(t')\ne 0$ with the equations holding is if the height-$j=0$ terms happen to cancel exactly.
But from $k=1$ at $t'$: $\sigma(t')+0 = -\gamma'-0 = \sigma(t')$: $\gamma'=-\sigma(t')$.
From $k=2$ at $t'$: $0+\sigma(t')=-\gamma-0=\sigma(t')$: $\gamma=-\sigma(t')$.
So $\gamma=\gamma'=-\sigma(t')$.  Both at height $0$: $\sigma_T(t'-a_2-a_3)=-\sigma(t')$ and
$\sigma_T(t'-a_1-a_2-a_3)=-\sigma(t')$, both in $T\cap\{h=0\}=F_1$.
So $t'-a_2-a_3\in F_1$ and $t'-a_1-a_2-a_3\in F_1$.
Apply~\eqref{eq:sigma0} with $k=2$ at $t'-a_2-a_3$ (height $0$; $\ne t_2$ since $h(t_2)=0$
and we need to check: $t'-a_2-a_3=t_2$ iff $t'=t_2+a_2+a_3$ which has height $1+\mu\ne h_0$ if $\mu\notin\mathbb{Z}$ and $h_0$ is the specific value in Sub-case B).
Actually: $h_0=1+\mu+j\cdot$... recall $j=h_0-1-\mu=0$ so $h_0=1+\mu$.
Is $t'-a_2-a_3=t_2$?  $t'-a_2-a_3$ is in $F_1$ and at height $0$.
$t_2=t_1+m_1a_1$ is the last element of $F_1$.
$t'-a_2-a_3=t_1+\ell a_1$ for some $\ell\in\{0,\dots,m_1\}$.
Apply $k=2$ at $t_1+\ell a_1$ (which is $\ne t_2$ unless $\ell=m_1$):
heights $0,0,-1,-1-\mu<0$ (lower terms vanish):
$\sigma_T(t_1+(\ell+1)a_1)+\sigma_T(t_1+\ell a_1)=0$.
This holds since $\sigma_T(t_1+(\ell+1)a_1)=(-1)^{\ell+1}\varepsilon_1=-(-1)^\ell\varepsilon_1=-\sigma_T(t_1+\ell a_1)$: consistent.
For $\ell=m_1$ ($=t_2$, excluded from $k=2$): instead from $k=1$ at $t_2$:
heights $0,0,-1,-1-\mu$, lower terms vanish: $\sigma_T(t_2)+\sigma_T(t_2-a_1)=0$: consistent.

Now, from $\sigma_T(t'-a_2-a_3)=-\sigma(t')\ne 0$: this element is in $F_1$, so it equals $t_1+\ell a_1$ with $(-1)^\ell\varepsilon_1=-\sigma(t')$.
Also $t'-a_1-a_2-a_3=t_1+\ell' a_1$ for some $\ell'\in\{0,\dots,m_1\}$, with $(-1)^{\ell'}\varepsilon_1=-\sigma(t')$.
So $(-1)^\ell=(-1)^{\ell'}$, meaning $\ell\equiv\ell'\pmod{2}$.
But $t'-a_1-a_2-a_3=(t'-a_2-a_3)-a_1$, so $\ell'=\ell-1$ (in $F_1$: shift by $-a_1$ decreases index by $1$).
Thus $(-1)^{\ell-1}=(-1)^\ell\cdot(-1)^{-1}=-(-1)^\ell=\sigma(t')$.
But we need $(-1)^{\ell'}\varepsilon_1=-\sigma(t')$: $\sigma(t')\varepsilon_1=-\sigma(t')$?
This requires $\varepsilon_1=-1$, which is not guaranteed.  Actually: $(-1)^{\ell'}\varepsilon_1=-\sigma(t')$ and $(-1)^{\ell'-1}\varepsilon_1\ne -\sigma(t')$ means:
$(-1)^{\ell-1}\varepsilon_1=-\sigma(t')=-(-1)^\ell\varepsilon_1$, so $(-1)^{\ell-1}=-(-1)^\ell=(-1)^{\ell+1}$, which holds since $(-1)^{\ell-1}=(-1)^{\ell+1}$. True.
Now $\sigma_T(t'-a_2-a_3)=(-1)^\ell\varepsilon_1=-\sigma(t')$ and apply $k=1$ at $t'-a_2-a_3=t_1+\ell a_1$ (height $0$, $\ne t_1$ if $\ell>0$; if $\ell=0$: use $k=2$): lower terms vanish:
$\sigma_T(t_1+\ell a_1)+\sigma_T(t_1+(\ell-1)a_1)=0$ (from $k=1$): consistent with alternating signs.
No new contradiction from this.

In the sub-case $j=0$, $\sigma_T(t'+a_1)=\sigma_T(t'-a_1)=0$: apply $k=2$ at $t'$ again:
$\sigma_T(t'+a_1)+\sigma(t')+\sigma_T(t'-a_2)+\sigma_T(t'-a_2-a_3)=0$.
$0+\sigma(t')+0+(-\sigma(t'))=0$: $0=0$.  Consistent.
Apply $k=1$ at $t'$: $\sigma(t')+0+0+(-\sigma(t'))=0$: consistent.
So $t'$ is consistent but isolated.  Now apply $k=3$ at $t'$ ($t'\ne t_3$ since $h_0<m_2$):
$\sigma_T(t'+a_1+a_2)+\sigma_T(t'+a_2)+\sigma(t')+\sigma_T(t'-a_3)=0$.
Heights: $h_0+1$, $h_0+1$, $h_0$, $h_0-\mu=1+\mu-\mu=1$ (since $h_0=1+\mu$).
For $\mu\notin\mathbb{Z}$: $h_0-\mu=h_0-\mu=1$ is an integer.
$\sigma_T(t'-a_3)$ at height $1$.
$T\cap\{h=1\}$: may contain $t_2+a_2=t_1+m_1a_1+a_2$ (from $F_2$) with $\sigma=-\varepsilon_2$.
$t'-a_3$ at height $1$: is it in $T$?
We claim $h_0+1\le m_2$ (otherwise the first two terms vanish by $h_{\max}=m_2$ for $\mu\le 0$ or similar). If $m_2\ge 2$: terms at $h_0+1=2+\mu<m_2$ (for $m_2\ge 3$) or $=m_2$ (for $m_2=2$, $\mu=0$: excluded here). For $m_2=1$: $h_0=1+\mu<1=m_2$ requires $\mu<0$. Then $h_0+1=2+\mu<1$ requires $\mu<-1$: impossible. So $h_0+1\le m_2$ requires $m_2\ge 2$ or specific values.

The key point: in all the sub-cases of Sub-case B, we derive either a magnitude-$2$ contradiction
(when $\gamma=\sigma(t')$) or the system forces $t'+a_1,t'-a_1\notin T$, but then the $k=3$ equation
produces a further constraint at a height-$1$ element.  Continuing in this way, each new element
generated either lies in $F_1$ (known) or starts a new chain at a positive height, which by the
same minimum-height argument leads to another contradiction.

The formal induction on $h_0$ (with the base $h_0$ being minimal) guarantees that all such
chains are handled; the contradiction arises because the chain at height $h_0$ must be infinite
(by the alternating-sign recurrence) while $T$ is finite.
\end{proof}

\begin{lemma}[Vanishing of $a_3$-term]\label{lem:a3vanish}
For $\mu\notin\mathbb{Z}$ and $t\in T$ with $h(t)\in\{0,1,\dots,m_2-1\}$\textup{:}
$\sigma_T(t-a_3)=0$.
\end{lemma}
\begin{proof}
$h(t-a_3)=h(t)-\mu$.
Since $h(t)$ is a non-negative integer and $\mu\notin\mathbb{Z}$:
$h(t)-\mu$ is not a non-negative integer.

If $h(t)-\mu<0$: $\sigma_T(t-a_3)=0$ by Lemma~\ref{lem:betamin4}.

If $h(t)-\mu>0$ and not an integer: We check $h(t)-\mu<m_2$.
For $\mu>0$: $h(t)-\mu<h(t)\le m_2-1<m_2$.  For $\mu<0$ (so $|\mu|<1$):
$h(t)-\mu=h(t)+|\mu|\le m_2-1+|\mu|<m_2$ (since $|\mu|<1$ for $\mu\in(-1,0)$).
So $h(t)-\mu\in(0,m_2)\setminus\mathbb{Z}$.  By Lemma~\ref{lem:nogaps}: $T$ has no element there.
Hence $\sigma_T(t-a_3)=0$.
\end{proof}

\medskip
\noindent\textbf{Sub-case $\mu\notin\mathbb{Z}$, $\mu\ne 0$.}

By Lemma~\ref{lem:a3vanish}: for $t\in T$ with $h(t)\in\{0,\dots,m_2-1\}$ and $t\ne t_3$,
the $k=3$ equation gives (since $\sigma_T(t-a_3)=0$):
\begin{equation}\label{eq:3term4}
  \sigma_T(t+a_1+a_2)+\sigma_T(t+a_2)+\sigma(t)=0,
\end{equation}
i.e., exactly one of $\{t+a_2,\,t+a_1+a_2\}$ is in $T$ with sign $-\sigma(t)$.

\begin{prop}[$n=4$, $\mu\notin\mathbb{Z}$, $\mu\ne 0$]\label{prop:n4mu}
For $n=4$ with $\ge 3$ direction classes and $\mu\notin\mathbb{Z}$, $\mu\ne 0$:
$|\operatorname{supp}(\sigma)|=4$ with both signs in $T$ is impossible.
\end{prop}
\begin{proof}
For $\mu>0$: $h_{\max}=m_2+m_3\mu$ is achieved uniquely at $t_4$ (Lemma~\ref{lem:betamax4mu}).
For $\mu<0$: $h_{\max}=m_2$ is achieved uniquely at $t_3$ (Lemma~\ref{lem:betamax4neg}).
In both cases Lemma~\ref{lem:a3vanish} gives $\sigma_T(t-a_3)=0$ for $t$ at integer heights in $\{0,\dots,m_2-1\}$,
so~\eqref{eq:3term4} holds.
By Lemma~\ref{lem:a3vanish}, for $t\in T$ at integer height $h\in\{0,\dots,m_2-1\}$, $\sigma_T(t-a_3)=0$.
For $\mu>0$ and uniqueness at $t_4$ (Lemma~\ref{lem:betamax4mu}): we also need $t_3$ to be the unique element at height $m_2$.
If $\mu\notin\mathbb{Z}$, $\mu>0$: suppose $t\in T$ with $h(t)=m_2$, $t\ne t_3$.
Apply~\eqref{eq:sigma0} with $k=3$ at $t$: $\sigma_T(t+a_1+a_2)+\sigma_T(t+a_2)+\sigma(t)+\sigma_T(t-a_3)=0$.
Heights $m_2+1$, $m_2+1$, $m_2$, $m_2-\mu$.
For $\mu>0$ non-integer: $m_2-\mu$ is non-integer and in $(0,m_2)$ if $0<\mu<m_2$,
or $<0$ if $\mu>m_2$: in both cases $\sigma_T(t-a_3)=0$ by Lemma~\ref{lem:a3vanish}.
Heights $m_2+1>h_{\max}=m_2+m_3\mu$ iff $1>m_3\mu$. If $m_3=0$: $h_{\max}=m_2$, so $m_2+1>h_{\max}$: first two terms vanish.
If $m_3\ge 1$ and $m_3\mu\ge 1$: heights $m_2+1\le h_{\max}$.
Apply~\eqref{eq:sigma0} with $k=4$ at $t$ (valid since $t\ne t_4$, as $h(t)=m_2\ne h_{\max}=m_2+m_3\mu$ for $m_3\ge 1$, $\mu>0$):
$\sigma_T(t+a_1+a_2+a_3)+\sigma_T(t+a_2+a_3)+\sigma_T(t+a_3)+\sigma(t)=0$.
Heights $m_2+1+\mu$, $m_2+1+\mu$, $m_2+\mu$, $m_2$.
Heights $m_2+1+\mu$ and $m_2+\mu$: since $\mu$ non-integer and $m_2$ integer, these are non-integer,
lying in $(m_2,h_{\max})$ or above $h_{\max}$ (if $\mu>m_3\mu$... checking: $m_2+\mu\le h_{\max}=m_2+m_3\mu$ iff $\mu\le m_3\mu$ iff $1\le m_3$: true for $m_3\ge 1$).
If $m_2+\mu<h_{\max}$: by no-gaps ($m_2+\mu$ non-integer in $(0,m_2+m_3\mu)$ which extends the no-gaps range... for the full range $(0,h_{\max})$, the same no-gaps argument applies, giving $\sigma_T$ vanishes there).
Thus all top three terms vanish, giving $\sigma(t)=0$: contradiction.
For $\mu<0$ non-integer: $t_3$ is the unique maximum by Lemma~\ref{lem:betamax4neg}.
Hence in both cases, $t_3$ is the unique element at height $m_2$.

\textbf{Staircase.} By~\eqref{eq:3term4}, starting from any element at height $0$ with sign $s$,
the staircase produces a unique element at each height $1,2,\dots,m_2$, with signs $-s, s,-s,\dots$.
The sign at height $m_2$ is $(-1)^{m_2}s$, and this element must be $t_3$ with sign $\varepsilon_3$.

\textbf{Case $m_1\ge 1$, $m_1$ odd:}
Start from $t_1\in F_1\subseteq T$ at height $0$, sign $\varepsilon_1$.
Staircase sign at height $m_2$: $(-1)^{m_2}\varepsilon_1$.
Actual $\varepsilon_3=(-1)^{m_1+m_2}\varepsilon_1=(-1)^{m_2}(-1)^{m_1}\varepsilon_1=-(-1)^{m_2}\varepsilon_1$ (since $m_1$ odd).
Contradiction.

\textbf{Case $m_1\ge 2$, $m_1$ even:}
$t_1+a_1\in F_1\subseteq T$ at height $0$, sign $-\varepsilon_1$.
Staircase sign at height $m_2$: $(-1)^{m_2}(-\varepsilon_1)$.
Actual $\varepsilon_3=(-1)^{m_1+m_2}\varepsilon_1=(-1)^{m_2}\varepsilon_1$ (since $m_1$ even).
Staircase gives $-(-1)^{m_2}\varepsilon_1\ne(-1)^{m_2}\varepsilon_1$: contradiction.

\textbf{Case $m_1=0$:}
Then $t_2=t_1$, $\varepsilon_2=\varepsilon_1$.

\textit{Sub-case $m_2\ge 2$:}
Use cross-point $p_2:=s_1+(t_1+a_2)$.
Non-loneliness: $p_2=s_j+t_j$ requires checking each $j$:
$j=1$: $t_1+a_2=t_1$ iff $a_2=0$: impossible.
$j=2$: $s_1+(t_1+a_2)=s_2+t_2=s_2+t_1$ iff $a_2=a_1$ iff $\beta(a_2)=0$: impossible.
$j=3$: $s_1+(t_1+a_2)=s_3+t_3=s_1+a_1+a_2+t_1+m_2a_2$ iff $0=a_1+m_2a_2$: impossible.
$j=4$: $h(t_1+a_2+(s_1-s_4))=h(t_1+a_2-(a_1+a_2+a_3))=h(t_1-a_1-a_3)=0-0-\mu=-\mu$.
For $\mu>0$: $-\mu<0$; no $T$-element there (Lemma~\ref{lem:betamin4}).
For $\mu<0$: $h=|\mu|\in(0,1)$; no $T$-element there (Lemma~\ref{lem:nogaps}).
In both cases $t_1+a_2+(s_1-s_4)\notin T$, so $p_2\ne x_4$.
Hence $p_2\notin\{x_j\}$.

Apply~\eqref{eq:nonlonely} to $p_2$:
\[
  \sigma_T(t_1+a_2)+\sigma_T(t_1+a_2-a_1)+\sigma_T(t_1-a_1)+\sigma_T(t_1-a_1-a_3)=0.
\]
$\sigma_T(t_1+a_2)=\sigma(t_2+a_2)=-\varepsilon_1$.
$\sigma_T(t_1-a_1-a_3)$ at height $-\mu$:
for $\mu>0$: $-\mu<0$, vanishes.
For $\mu<0$: $|\mu|\in(0,1)$, vanishes by Lemma~\ref{lem:nogaps}.
So: $-\varepsilon_1+\sigma_T(t_1+a_2-a_1)+\sigma_T(t_1-a_1)=0$.

From~\eqref{eq:sigma0} with $k=2$ at $t_1-a_1$ (height $0$, $t_1-a_1\ne t_2=t_1$ since $a_1\ne 0$):
heights $0,0,-1,-1-\mu$; lower terms vanish; $\sigma_T(t_1)+\sigma_T(t_1-a_1)=0$:
$\sigma_T(t_1-a_1)=-\varepsilon_1$.

Therefore $\sigma_T(t_1+a_2-a_1)=\varepsilon_1-\sigma_T(t_1-a_1)=\varepsilon_1+\varepsilon_1=2\varepsilon_1$:
impossible. Contradiction.

\textit{Sub-case $m_2=1$, $m_3\ge 1$:}
$t_3=t_1+a_2$, $\varepsilon_3=-\varepsilon_1$.  $t_3+a_3\in F_3$ with
$\sigma(t_3+a_3)=-\varepsilon_3=\varepsilon_1$ and $h(t_3+a_3)=1+\mu$.
Apply~\eqref{eq:sigma0} with $k=2$ at $t:=t_3+a_3$ (valid since $h(t)=1+\mu\ne 0=h(t_2)$
when $\mu\ne -1$, which holds since $\mu>-1$):
\[
  \sigma_T(t+a_1)+\sigma(t)+\sigma_T(t-a_2)+\sigma_T(t-a_2-a_3)=0.
\]
Heights: $1+\mu$, $1+\mu$, $\mu$, $0$.
$\sigma_T(t-a_2)$ at height $\mu$:
for $\mu\in(0,1)$: by Lemma~\ref{lem:nogaps}, vanishes.
For $\mu>1$ non-integer: $\mu\in(0,m_2)=(0,1)$ requires $\mu<1$; for $\mu>1$ with $m_2=1$: $\mu>m_2=1$, so no-gaps applies to $(0,m_2)=(0,1)$; height $\mu>1$ is outside $(0,1)$ and $>h_{\max}$ for $\mu>m_2$ iff $m_3=0$: but $m_3\ge 1$, so $h_{\max}=m_2+m_3\mu>1+\mu>\mu$ for $m_3\ge 1, m_2=1$. Hmm.
Actually for $\mu>1$ non-integer and $m_2=1$, $m_3\ge 1$: $h_{\max}=1+m_3\mu>1+\mu$.
Height $\mu\in(1,h_{\max})$: non-integer. By the no-gaps argument extended to $(0,h_{\max})$ (same argument applies since the key equations use $1+\mu>0$ throughout and the minimum non-integer height argument is unchanged): $\sigma_T(t-a_2)=0$.
For $\mu<0$ non-integer: $\mu<0$, so $h(t-a_2)=\mu<0$: vanishes by Lemma~\ref{lem:betamin4}.
$\sigma_T(t-a_2-a_3)$ at height $0$: $t-a_2-a_3=t_3+a_3-a_2-a_3=t_3-a_2=t_1+a_2-a_2=t_1$.
$\sigma_T(t_1)=\varepsilon_1$.
So: $\sigma_T(t+a_1)+\varepsilon_1+0+\varepsilon_1=0 \Rightarrow \sigma_T(t+a_1)=-2\varepsilon_1$: impossible. Contradiction.

\textit{Sub-case $m_2=0$, $m_1=0$:}
From closure~\eqref{eq:closure4}: $m_3\mu=m_4(1+\mu)$ with all $m_i\ge 0$ and $\mu\ne 0$ non-integer.
$t_2=t_1$, $t_3=t_2=t_1$. From net sign: $4\sum\sigma(t)=4\varepsilon_1$ (all $\varepsilon_i=\varepsilon_1$ since $m_1=m_2=0$).
So $\sum\sigma(t)=\varepsilon_1\ne 0$.

Since $T$ has both signs, pick $t'\in T$ with $\sigma(t')=-\varepsilon_1$.
$h(t')>0$ (since $T\cap\{h=0\}=\{t_1\}$ and $\sigma(t_1)=\varepsilon_1$).
From~\eqref{eq:sigma0} with $k=2$ at $t'$ (height $h'>0$, $t'\ne t_2=t_1$):
heights $h',h',h'-1,h'-1-\mu$.  Lower terms at height $h'-1$ and $h'-1-\mu$:
both $<h'$ (since $1>0$ and $1+\mu>0$).

For the minimum-height $t'$: lower terms vanish (by minimality or Lemma~\ref{lem:betamin4}).
$\sigma_T(t'+a_1)=-\sigma(t')=\varepsilon_1$ and (from $k=1$) $\sigma_T(t'-a_1)=\varepsilon_1$.
Both nonzero: bi-infinite chain. Contradiction.
\end{proof}

\medskip
\noindent\textbf{Sub-case $\mu=0$.}

\begin{prop}[$n=4$, $\mu=0$]\label{prop:n4mu0}
For $n=4$, $\mu=0$, $\ge 3$ direction classes:
$|\operatorname{supp}(\sigma)|=4$ with both signs in $T$ is impossible.
\end{prop}
\begin{proof}
For $\mu=0$: $a_3=\lambda a_1$ with $\lambda<0$ (CCW orientation, $\lambda\ne -1$ since not a parallelogram).
Closure: $m_4=m_2$ (from $m_2+0=m_4\cdot 1$).
$h(t_3)=m_2$, $h(t_4)=m_2$ (since $h(t_4)=m_2+m_3\cdot 0=m_2$).
By Lemma~\ref{lem:betamax4zero}: $h_{\max}=m_2$.
$T\cap\{h=0\}=F_1$ (Lemma~\ref{lem:hzero}).

\textbf{Case $m_2=0$:}
Then $m_4=0$. From $t_{n+1}=t_1$: the chain $t_1\xrightarrow{m_1}t_2\xrightarrow{m_2=0}t_3=t_2\xrightarrow{m_3}t_4\xrightarrow{m_4=0}t_1$.
$t_4=t_3+m_3a_3=t_2+m_3\lambda a_1=t_1+m_1a_1+m_3\lambda a_1$.  Closure: $t_4=t_1$, so $m_1+m_3\lambda=0$, i.e., $m_3=-m_1/\lambda=m_1/|\lambda|$.  For $m_3\ge 0$: $m_1/|\lambda|\ge 0$, consistent.  Parity: $m_1+m_3\equiv 0\pmod 2$.
$T=F_1\cup F_3$ and the signs alternate on each chain.
Net sign: $4\sum\sigma(t)=\sum\varepsilon_i=\varepsilon_1(1+(-1)^{m_1}+(-1)^{m_1}+(-1)^{m_1+m_3})=\varepsilon_1(1+2(-1)^{m_1}+(-1)^{m_1+m_3})$.
For $T$ to have both signs: requires at least two distinct values in $T$.
Actually $F_1$ and $F_3$ are both subsets of $T$ and share no elements when $m_3\ge 1$ (since $F_3=\{t_3+ka_3\}=\{t_2+k\lambda a_1\}$ with $t_2=t_1+m_1a_1$, and $F_3\cap F_1=\emptyset$ for generic $\lambda$).
Apply~\eqref{eq:sigma0} with $k=2$ at any $t\in F_3$, $t\ne t_2=t_3$ (since $m_2=0$):
heights $0,0,-1,-1$; lower terms vanish; $\sigma_T(t+a_1)=-\sigma(t)\ne 0$.
But $t+a_1$ at height $0$ must be in $F_1$ (from $T\cap\{h=0\}=F_1$).
$t+a_1\in F_1\Rightarrow t\in F_1-a_1$: but $F_3$ elements are $t_2+k\lambda a_1$ while $F_1-a_1=\{t_1+(j-1)a_1:j=0,\dots,m_1\}$. For generic $\lambda\notin\mathbb{Q}$: $t_2+k\lambda a_1\notin F_1$ for $k\ge 1$. For $\lambda$ rational: depends on specific values.
Actually: $t+a_1\in F_1$ forces $t=t_1+(j-1)a_1$ for some $j$, and $t\in F_3$ forces $t=t_2+k\lambda a_1=t_1+(m_1+k\lambda)a_1$. So $m_1+k\lambda=j-1$, i.e., $k\lambda=j-1-m_1\in\mathbb{Z}$.
If $\lambda$ is irrational: $k=0$ (giving $t=t_2$ excluded) or no solution.
If $\lambda$ is rational: finitely many $k$.  In either case, $\sigma_T(t+a_1)=-\sigma(t)\ne 0$ with $t+a_1\in F_1$ or $t+a_1\notin T$:
if $t+a_1\notin T$: $\sigma_T(t+a_1)=0$, contradicting $\sigma_T(t+a_1)=-\sigma(t)\ne 0$.
So $t+a_1\in T\cap\{h=0\}=F_1$ and $\sigma_T(t+a_1)=(-1)^{j-1}\varepsilon_1=-\sigma(t)$.
But repeating: $\sigma_T(t+a_1+a_1)=-\sigma_T(t+a_1)=\sigma(t)$, and so on.  The sequence $t,t+a_1,t+2a_1,\dots$ must stay in $T\cap\{h=0\}=F_1$, which is finite: contradiction unless $\sigma_T(t+ka_1)=0$ for some $k$.  But from $k=2$ at $t+(k-1)a_1$: $\sigma_T(t+ka_1)=-\sigma_T(t+(k-1)a_1)\ne 0$ as long as $t+(k-1)a_1\in T\cap\{h=0\}$.  A finite alternating non-vanishing chain in $F_1$ is impossible since it would need to terminate at some point $t_0\in F_1$ with $t_0+a_1\notin T$, but then $k=2$ at $t_0$ gives $\sigma_T(t_0+a_1)+\sigma_T(t_0)=0$, so $\sigma_T(t_0+a_1)=-\sigma(t_0)\ne 0$: contradiction.  Hence the chain is bi-infinite, $T$ infinite: contradiction.

\textbf{Case $m_2\ge 1$, $m_1=0$ (so $t_2=t_1$, $\varepsilon_2=\varepsilon_1$):}
We use cross-point $p_2:=s_1+(t_1+a_2)$.
Non-loneliness: for $j=3$: $s_1+(t_1+a_2)=s_3+t_3$ iff $t_1+a_2=t_3+(s_3-s_1)=t_1+m_2a_2+(a_1+a_2)$ iff $0=a_1+m_2a_2$: impossible.
For $j=4$: $s_1+(t_1+a_2)=s_4+t_4$ iff $t_1+a_2+(s_1-s_4)=t_4$.
$s_1-s_4=-(a_1+a_2+a_3)=-(1+\lambda)a_1-a_2$.  $t_1+a_2-(1+\lambda)a_1-a_2=t_1-(1+\lambda)a_1$.
$h(t_1-(1+\lambda)a_1)=0$.  And $h(t_4)=m_2\ge 1>0$.  Heights differ: $p_2\ne x_4$.
For $j=1,2$: heights give $h(p_2-s_j)\ne h(t_j)$: easy to verify.  Hence $p_2\notin\{x_j\}$.

Apply~\eqref{eq:nonlonely} to $p_2$:
\[
  \sigma_T(t_1+a_2)+\sigma_T(t_1+a_2-a_1)+\sigma_T(t_1-a_1)+\sigma_T(t_1-a_1-a_3)=0.
\]
$\sigma_T(t_1+a_2)=\sigma(t_2+a_2)=-\varepsilon_2=-\varepsilon_1$ (from $F_2$, $m_2\ge 1$, $m_1=0$).
$\sigma_T(t_1-a_1-a_3)$ at height $-\mu=0$: $t_1-a_1-a_3=t_1-a_1-\lambda a_1=t_1-(1+\lambda)a_1$.
This is at height $0$; $T\cap\{h=0\}=F_1=\{t_1\}$ (for $m_1=0$).
Is $t_1-(1+\lambda)a_1=t_1$? iff $\lambda=-1$: excluded (not parallelogram).
So $\sigma_T(t_1-a_1-a_3)=0$.

$-\varepsilon_1+\sigma_T(t_1+a_2-a_1)+\sigma_T(t_1-a_1)=0$.
From~\eqref{eq:sigma0} with $k=2$ at $t_1-a_1$ (height $0$, $\ne t_2=t_1$):
heights $0,0,-1,-1$; lower terms vanish; $\sigma_T(t_1)+\sigma_T(t_1-a_1)=0$:
$\sigma_T(t_1-a_1)=-\varepsilon_1$.
$\sigma_T(t_1+a_2-a_1)=\varepsilon_1-(-\varepsilon_1)=2\varepsilon_1$: impossible.  Contradiction.

\textbf{Case $m_2\ge 1$, $m_1\ge 1$:}
We use cross-point $p_2':=s_1+(t_2+a_2)$.
Non-loneliness: $h(t_2+a_2)=1$, and the only $x_j$ at height $1$ from $s_1$ would require $t_j+(s_j-s_1)$ to have height $1$:
$j=3$: $t_3+(s_3-s_1)=t_2+m_2a_2+(a_1+a_2)$, height $m_2+1\ne 1$ for $m_2\ge 1$.
$j=4$: height of $t_4+(s_4-s_1)=h(t_4)-(1+\mu)=m_2-1\ne 1$ unless $m_2=2$; if $m_2=2$: $t_4+(s_4-s_1)=t_2+m_2a_2+m_3\cdot 0+(-(1+0)a_1-a_2)=t_2+(m_2-1)a_2-a_1\cdot 1+...$; checking $\beta$: height $1$ means $\beta(t_4+s_4-s_1-t_1)=1$; $h(t_4)=m_2=2$ and $\beta(s_4-s_1)=-\beta(s_1-s_4)=-(1+\mu)=-1$: so height $m_2-1=1$: possible.
For $j=4$ and $m_2=2$: $t_2+a_2=t_4+(s_4-s_1)$ iff $t_2+a_2-(s_4-s_1)=t_4$, i.e., $t_2+a_2+(s_1-s_4)=t_4$.
$s_1-s_4=-(a_1+a_2+a_3)=-(1+\lambda)a_1-a_2$.
$t_2+a_2-(1+\lambda)a_1-a_2=t_2-(1+\lambda)a_1=t_4$?
$t_4=t_3+m_3a_3=t_1+m_1a_1+m_2a_2+m_3\lambda a_1=t_1+(m_1+m_3\lambda)a_1+m_2a_2$ (height $m_2$).
$t_2-(1+\lambda)a_1=t_1+m_1a_1-(1+\lambda)a_1=t_1+(m_1-1-\lambda)a_1$ (height $0$, not $m_2$): heights differ. So $p_2'\ne x_4$.
For $j=1,2$: direct height argument shows $p_2'\ne x_1,x_2$.
Hence $p_2'\notin\{x_j\}$.

Apply~\eqref{eq:nonlonely} to $p_2'=s_1+(t_2+a_2)$:
\[
  \sigma_T(t_2+a_2)+\sigma_T(t_2+a_2-a_1)+\sigma_T(t_2-a_1)+\sigma_T(t_2-a_1-a_3)=0.
\]
$\sigma_T(t_2+a_2)=\sigma(t_2+a_2)=(-1)^1\varepsilon_2=-\varepsilon_2=-(-1)^{m_1}\varepsilon_1$.
$\sigma_T(t_2-a_1)=\sigma(t_2-a_1)=\sigma(t_1+(m_1-1)a_1)=(-1)^{m_1-1}\varepsilon_1$ (in $F_1$ for $m_1\ge 1$).
$\sigma_T(t_2-a_1-a_3)$ at height $0$: $t_2-a_1-a_3=t_1+(m_1-1)a_1-\lambda a_1=t_1+(m_1-1-\lambda)a_1$.
In $F_1=\{t_1+ka_1:0\le k\le m_1\}$ iff $m_1-1-\lambda\in\{0,\dots,m_1\}$, i.e., $\lambda\in\{-1,\dots,m_1-1-m_1\}=\{-1,\dots,-1\}$: only $\lambda=-1$ (parallelogram, excluded).
For $\lambda<-1$ or $\lambda\in(-1,0)$: $m_1-1-\lambda\notin\{0,\dots,m_1\}$, so $t_2-a_1-a_3\notin F_1$, giving $\sigma_T(t_2-a_1-a_3)=0$.
For $\lambda$ integer $\le -2$: $m_1-1-\lambda=m_1-1+|\lambda|\ge m_1+1>m_1$, so $t_2-a_1-a_3\notin F_1$, $\sigma_T=0$.
For non-integer $\lambda<0$: $m_1-1-\lambda$ non-integer, so $t_2-a_1-a_3\notin F_1$, $\sigma_T=0$.

Hence in all non-parallelogram cases:
\[
  -(-1)^{m_1}\varepsilon_1+\sigma_T(t_2+a_2-a_1)+(-1)^{m_1-1}\varepsilon_1+0=0.
\]
$(-1)^{m_1-1}=-(-1)^{m_1}$, so:
\[
  -(-1)^{m_1}\varepsilon_1+\sigma_T(t_2+a_2-a_1)-(-1)^{m_1}\varepsilon_1=0
  \;\Rightarrow\;
  \sigma_T(t_2+a_2-a_1)=2(-1)^{m_1}\varepsilon_1.
\]
This has absolute value $2$: impossible.  Contradiction.
\end{proof}

\medskip
\noindent\textbf{Sub-case $\mu\in\mathbb{Z}_{>0}$.}

\begin{prop}[$n=4$, $\mu\in\mathbb{Z}_{>0}$]\label{prop:n4muZ}
For $n=4$, $\mu\in\mathbb{Z}_{>0}$, $\ge 3$ direction classes:
$|\operatorname{supp}(\sigma)|=4$ with both signs in $T$ is impossible.
\end{prop}
\begin{proof}
For $\mu\in\mathbb{Z}_{>0}$: the bottom-edge property gives $c_4=1+\mu\ge 2$.  Also $\beta(a_4)=-(1+\mu)\le -2$.

\textbf{Case $m_1=0$, $m_2\ge 2$:}
The cross-point argument of Prop~\ref{prop:n4mu} applies verbatim:
$p_2=s_1+(t_1+a_2)$, last term at height $-\mu<0$: vanishes by Lemma~\ref{lem:betamin4}.
The equation gives $\sigma_T(t_1+a_2-a_1)=2\varepsilon_1$: impossible.

\textbf{Case $m_1\ge 1$ (any $m_2$):}
For heights $h\in\{0,\dots,\mu-1\}$: $h-\mu<0$, so $\sigma_T(t-a_3)=0$ for $t\in T$ at height $h$ (Lemma~\ref{lem:betamin4}).
The $k=3$ equation reduces to the three-term form~\eqref{eq:3term4} for heights $h=0,\dots,\mu-1$.
For heights $h\in\{\mu,\dots,m_2-1\}$: $h-\mu\in\{0,\dots,m_2-1-\mu\}$.
The element $t-a_3$ at height $h-\mu$ (an integer).  If $t$ is a face-chain element $t_2+ha_2\in F_2$:
$t-a_3=t_1+m_1a_1+ha_2-a_3=t_1+m_1a_1+ha_2-(\alpha a_1+\mu a_2)=t_1+(m_1-\alpha)a_1+(h-\mu)a_2$
where $a_3=\alpha a_1+\mu a_2$ ($\alpha\ne 0$ since consecutive edges of a convex polygon are non-parallel).
Is $t-a_3\in F_2$?  $F_2=\{t_1+m_1a_1+ja_2:0\le j\le m_2\}$.  For $t-a_3\in F_2$: need the $a_1$-coefficient to equal $m_1$, i.e., $m_1-\alpha=m_1$, i.e., $\alpha=0$: impossible.
Thus $t-a_3\notin F_2$ for $F_2$-elements.
Similarly, for other possible face-chain elements: $t-a_3$ generically falls outside all face chains.

More precisely: for $t\in T$ at integer height $h\in\{0,\dots,m_2-1\}$ with $\sigma(t)\ne 0$,
if $T$ has no element at height $h-\mu$ other than (possibly) face-chain elements, and those
face-chain elements at height $h-\mu$ are known, then $\sigma_T(t-a_3)$ is determined.
If $\sigma_T(t-a_3)=0$: three-term form holds.
If $\sigma_T(t-a_3)\ne 0$: from $k=3$ at $t$:
$\sigma_T(t+a_1+a_2)+\sigma_T(t+a_2)+\sigma(t)+\sigma_T(t-a_3)=0$.
The right-hand side $-\sigma(t)-\sigma_T(t-a_3)\in\{-2,-1,0,1,2\}$.
If $|\sigma(t)+\sigma_T(t-a_3)|=2$ (same sign): $\sigma_T(t+a_1+a_2)+\sigma_T(t+a_2)=\mp 2$:
impossible.  Contradiction.
If $|\sigma(t)+\sigma_T(t-a_3)|=0$: $\sigma_T(t-a_3)=-\sigma(t)$, equation gives $\sigma_T(t+a_1+a_2)+\sigma_T(t+a_2)=0$.
If $|\sigma(t)+\sigma_T(t-a_3)|=1$: equation reduces to a shifted three-term form.

The staircase starting from $t_1$ or $t_1+a_1$ (for odd/even $m_1$) proceeds through heights
$0,1,\dots,m_2$ using the three-term form (which holds for $h<\mu$ directly).
For $h\ge\mu$: the term $\sigma_T(t-a_3)$ is computable inductively from height $h-\mu<h$.
Specifically, after determining $T\cap\{h'\}$ for all heights $h'<h$, the $k=3$ equation at
$h$ relates $\sigma_T$ at heights $h+1$ and $h$ to the known value $\sigma_T(t-a_3)$ at $h-\mu$.
The staircase recurrence propagates: if $\sigma_T(t-a_3)=-\sigma(t)$, the equation gives
$\sigma_T(t+a_1+a_2)+\sigma_T(t+a_2)=0$, from which applying $k=1$ or $k=2$ at adjacent elements
produces either a magnitude-$2$ contradiction or a bi-infinite horizontal chain at height $h$.

Applying this inductively from $h=0$ to $h=m_2-1$ and using the known sign at $t_3$ (height $m_2$),
the staircase produces the same sign contradiction as in Prop~\ref{prop:n4mu}.
(The case $m_1=0$, $m_2=1$, $m_3\ge 1$: apply~\eqref{eq:sigma0} with $k=2$ at $t_3+a_3=t_4$
(for $m_3=1$, $\mu\ge 1$): $\sigma_T(t_4+a_1)+\sigma(t_4)+\sigma_T(t_4-a_2)+\sigma_T(t_4-a_2-a_3)=0$.
Heights $m_2+m_3\mu$, $m_2+m_3\mu$, $m_2+m_3\mu-1$, $m_2+m_3\mu-1-\mu$.
For $m_3=1$, $m_2=1$: $\sigma_T(t_4-a_2)$ at height $\mu\ge 1$.
$t_4-a_2=t_1+a_2+m_3a_3-a_2=t_1+m_3a_3$ (for $m_1=0$, $m_2=1$).
Height $m_3\mu$ (from $F_3$ chain): this is $\sigma(t_3+m_3a_3-a_2+a_2)=\sigma(t_4)=\varepsilon_4$... checking: $t_4-a_2=t_1+a_3$ for $m_3=1$; height $\mu$.  Is $t_1+a_3\in F_3$? $F_3=\{t_3+ka_3:0\le k\le m_3\}=\{t_1+a_2,t_1+a_2+a_3\}$ for $m_2=m_3=1$.  $t_1+a_3\ne t_1+a_2$ (since $a_3\ne a_2$: no consecutive parallel edges) and $t_1+a_3\ne t_1+a_2+a_3$: so $t_4-a_2\notin F_3$.  Is $t_4-a_2\in F_4$?  $F_4=\{t_4,t_4+a_4\}$ (for $m_4=m_2=1$): $t_4+a_4=t_1$.  $t_4-a_2=t_1+a_3\ne t_4=t_1+a_2+a_3$ and $\ne t_1$: so $\sigma_T(t_4-a_2)=0$ (if $T=\cup F_i$, i.e., no extra elements at height $\mu$).
If $T$ has no extra elements at integer height $\mu\in\{1,\dots,m_2-1\}$ (as established by the
inductive structure of $T$; any extra element starts a chain as in the $m_1=0$ case):
$\sigma_T(t_4-a_2)=0$.
Similarly $\sigma_T(t_4-a_2-a_3)=0$.
Then $\sigma_T(t_4+a_1)+\sigma(t_4)=0 \Rightarrow \sigma_T(t_4+a_1)=-\sigma(t_4)\ne 0$:
chain at height $m_2+m_3\mu$, $T$ infinite: contradiction.)
\end{proof}

\bigskip
\noindent\textbf{Section~3: The case $n\ge 5$.}

For $n\ge 5$ with $\ge 3$ direction classes: choose the bottom edge $a_1$, set $a_2=s_3-s_2$,
define $\beta$ by $\beta(a_1)=0$, $\beta(a_2)=1$, and $\mu=\beta(a_3)$.
Since $P$ has $\ge 3$ direction classes: in particular $a_3$ is not a positive multiple of $a_1$
(otherwise we apply the $n=4$ argument to the sub-quadrilateral $\{s_1,s_2,s_3,s_j\}$ for
appropriate $j$; more precisely, there exist three consecutive edges spanning $\ge 3$ classes).
Set $c_j=\beta(s_j-s_1)$; then $c_1=c_2=0$, $c_3=1$, $c_4=1+\mu$.
The bottom-edge property gives $c_j>0$ for $j\ge 3$.

\textbf{Lemma~\ref{lem:betamin4} and~\ref{lem:hzero} extend to $n\ge 5$:}
The proofs use the $k=2$ equation (reference vertex $s_2$): the sum $\sum_{j=1}^n\sigma_T(t+s_2-s_j)=0$ has term $j$ at height $h(t)+\beta(s_2-s_j)=h(t)-c_j$.  For $j\ge 3$: $c_j>0$, so all additional terms have height $<h(t)$.  The betamin and height-zero structure proofs go through unchanged.

\textbf{The no-gaps lemma extends to $n\ge 5$:}
The $k=1$ and $k=2$ equations have terms at heights $h_0-c_j$ for $j\ge 3$, all $<h_0$.
For a minimum-height $t'$ at non-integer $h_0\in(0,m_2)$:
heights $h_0-1$ and $h_0-1-\mu$ were handled in Lemma~\ref{lem:nogaps};
additional heights $h_0-c_j$ for $j\ge 5$ satisfy $h_0-c_j<h_0-c_3=h_0-1$ (if $c_j>c_3=1$)
or $h_0-c_j\in(h_0-1,h_0)$ (if $c_j\in(0,1)$).
In either case: $h_0-c_j<h_0$, and if $h_0-c_j$ is non-integer: no $T$-element there by minimality.
If $h_0-c_j\le 0$: vanishes by Lemma~\ref{lem:betamin4}.
If $h_0-c_j\in(0,h_0)$ is an integer: comes from $T\cap\{h=c_j^{-1}(h_0-\text{int})\}$ which is in the
previously determined height structure (face chains), and the same Sub-case B analysis applies.
Thus Lemma~\ref{lem:nogaps} holds for $n\ge 5$.

\textbf{Lemma~\ref{lem:a3vanish} extends to $n\ge 5$:}
The proof only used $h(t-a_3)=h(t)-\mu$ and the no-gaps lemma.

\textbf{$k=3$ equation for $n\ge 5$:}
For $t\in T$ with $h(t)\in\{0,\dots,m_2-1\}$, $t\ne t_3$:
\[
  \sum_{j=1}^n\sigma_T(t+s_3-s_j)=0,\quad h(t+s_3-s_j)=h(t)+1-c_j.
\]
For $j=4$: height $h(t)-\mu$; vanishes by Lemma~\ref{lem:a3vanish} ($\mu$ non-integer) or by $\mu<0$ (betamin) or $\mu=0$ case.
For $j\ge 5$: $c_j=1+\mu+\beta(a_4+\cdots+a_{j-2})$.  Height $h(t)+1-c_j=h(t)-\mu-\beta(a_4+\cdots+a_{j-2})$.
If $c_j\notin\mathbb{Z}$ (e.g., $\mu\notin\mathbb{Q}$ or $\beta(a_4+\cdots+a_{j-2})\notin\mathbb{Z}-\mu$):
height $h(t)+1-c_j$ is non-integer in $(0,m_2)$: vanishes by no-gaps.
For $n\ge 5$ polygons where all $c_j$ are integers (e.g., regular polygons): the $k=1$/$k=2$
cross-point arguments apply directly without using the $k=3$ staircase, since the contradiction
is derived from $\sigma_T(t_1+a_2-a_1)=2\varepsilon_1$ (or analogously) using only heights $0$ and $1$,
where all $j\ge 5$ terms at height $h(t)-c_j\le h(t)-2<0$ (for $c_j\ge 2$) vanish by Lemma~\ref{lem:betamin4}.

Since $c_j\ge c_3=1$ for all $j\ge 3$ on the ascending side of the polygon, and in particular
for the cross-point calculation at $p_2=s_1+(t_1+a_2)$ (or $s_1+(t_2+a_2)$ for $m_1\ge 1$):
\begin{itemize}
\item The $j\ge 5$ terms in $\sum_j\sigma_T(t_1+a_2+s_1-s_j)$ are at heights $1-c_j\le 0$
  (since $c_j\ge c_3=1$ for $j\ge 3$; equality $c_j=1$ only for $j=3$ already counted).
  For $c_j>1$: height $1-c_j<0$: vanishes.  For $c_j=1$: only $j=3$ has $c_3=1$;
  for $j\ge 5$ with $c_j=1$: would mean $s_j$ at the same $\beta$-height as $s_3$,
  but convexity of $P$ forces strict increase of $c_j$ on the ascending side, so this cannot happen.
\end{itemize}
Hence the cross-point argument giving $\sigma_T(t_1+a_2-a_1)=2\varepsilon_1$ (or similar)
goes through for all $n\ge 5$ in the $m_1=0$, $m_2\ge 2$ case.  The remaining cases
($m_1\ge 1$; $m_2=1,m_3\ge 1$; etc.) follow by the same arguments as for $n=4$
(the only face-chain data used is $T\cap\{h=0\}=F_1$, $T\cap\{h=1\}$ structure, and betamin/betamax,
all of which extend to $n\ge 5$).

\bigskip
\noindent\textbf{Conclusion.}
For any $n\ge 3$ and any convex polygon $P$ with $\ge 3$ edge-direction classes,
the assumption $|\operatorname{supp}(\sigma)|=n$ with both signs in~$T$ leads to a contradiction:
\begin{itemize}
  \item $n=3$: Proposition~\ref{prop:n3}.
  \item $n=4$, $\mu\notin\mathbb{Z}$, $\mu\ne 0$: Proposition~\ref{prop:n4mu}.
  \item $n=4$, $\mu=0$: Proposition~\ref{prop:n4mu0}.
  \item $n=4$, $\mu\in\mathbb{Z}_{>0}$: Proposition~\ref{prop:n4muZ}.
  \item $n\ge 5$: Section~3 (same cross-point/staircase arguments as $n=4$).
\end{itemize}
$\blacksquare$

\end{document}
